\documentclass[10pt,a4paper]{amsart}
\usepackage[margin=19mm]{geometry}
\usepackage{amsmath,amssymb,mathtools}
\usepackage[scr=rsfs,cal=euler]{mathalfa}
\usepackage{xfrac}
\usepackage[unicode,hidelinks,bookmarks=false]{hyperref}
\newcommand{\ie}{i.e.\ }
\newcommand{\cf}{cf.\ }
\newcommand{\resp}{respectively\ }
\newcommand{\Subsection}{\subsection}
\providecommand{\nobreakdash}{\nobreak\hbox{-}}
\setlength{\emergencystretch}{2em}
\multlinegap0pt
\usepackage{amssymb}
\usepackage{mathtools}
\usepackage{enumerate}
\usepackage{xcolor}
\newtheorem{theorem}{Theorem}
\newcommand{\AI}{A\upI}
\newcommand{\AR}{A\upR}
\newcommand{\I}{\mathrm{i}}    % imaginary unit
\newcommand{\LambdaR}{\Lambda\uptR}
\newcommand{\di}{\mathrm{d}}   % differential
\newcommand{\imdel}{\mathop{\textup{Im}}}
\newcommand{\inner}[1]{\langle#1\rangle} 
\newcommand{\redel}{\mathop{\textup{Re}}}
\newcommand{\upI}{^{\textup{I}}}
\newcommand{\upR}{^{\textup{R}}}
\newcommand{\uptR}{^{\!\textup{R}}}
\title{Reconstruction of non-self-adjoint anisotropic and complex inclusions in the Calder\'on problem}
\author{Henrik Garde, David Johansson, Thanasis Zacharopoulos}
\date{Source: arXiv:2605.05021v1 (CC BY 4.0)}
\begin{document}
\maketitle

\noindent\emph{Source and licence.} Reconstruction of non-self-adjoint anisotropic and complex inclusions in the Calder\'on problem. Version arXiv:2605.05021v1 (CC BY 4.0), distributed by arXiv under the Creative Commons Attribution 4.0 International licence, http://creativecommons.org/licenses/by/4.0/, as recorded in the arXiv licence metadata for this exact version and shown on the abstract page https://arxiv.org/abs/2605.05021v1. Full licence notice: \texttt{licenses/arxiv-2605.05021-CC-BY-4.0.txt}.\par\smallskip

\noindent\textit{Editorial scope.} Counted unit: the article's theorem thm:improvedmono (source0.tex lines 930--937). The article's complete original proof of it is delivered with it (source0.tex lines 939--985, one \texttt{proof} environment of 47 lines). No mathematics is written, repaired or re-derived: the statement and the proof are the article's own lines, in the article's own notation, and no retained line is substituted or re-flowed.  No further source-local context is needed: every cross-reference of every retained line resolves inside the two delivered spans.  Nothing was omitted from any retained span, so the package carries no omission record.  No retained line contains a citation, so the wrapper carries no bibliography and no bibliography entry.  No retained line contains a figure environment, an image-inclusion command or drawing code, so no image is delivered and the package declares no asset dependency.  Editorial formatting only, in addition: an AMS \texttt{amsart} preamble that restates the article's own declarations and macros used by the retained lines, namely the environments \texttt{theorem} on separate counters, together with the standard packages those lines need.  The retained environments print in this document as Theorem 1 in order of appearance, whereas the article prints the same items with the numbering of its own full document.  The complete native source of the article as distributed in the arXiv e-print for this exact version is bundled unchanged as \texttt{sources/199915/source0.tex}.
	\begin{theorem}\label{thm:improvedmono}
		Let $A_1,A_2\in\mathcal{H}(\Omega)$ and $f\in L^2_\diamond(\Gamma)$. For $j\in\{1,2\}$ we denote $\Lambda_j = \Lambda(A_j)$ and $u_j = u_f^{A_j}$. Then
		%
		\begin{align*}
			\inner{f,(\LambdaR_1-\LambdaR_2) f} &\geq \int_\Omega (\AR_2-\AR_1)\nabla u_2\cdot \overline{\nabla u_2}\,\di x + 2\imdel\int_\Omega \AI_1\nabla u_1\cdot\overline{\nabla u_2}\,\di x,  \\
			\inner{f,(\LambdaR_1-\LambdaR_2) f} &\leq \int_\Omega \AR_2(\AR_1)^{-1}(\AR_2-\AR_1)\nabla u_2\cdot\overline{\nabla u_2}\,\di x + 2\imdel \int_\Omega \AI_2\nabla u_1\cdot\overline{\nabla u_2}\,\di x.
		\end{align*}
	\end{theorem}

	\begin{proof}
		We start by noting that 
		%
		\begin{align} \label{eq:mixed1}
			\redel \int_\Omega \AR_1\nabla u_1\cdot\overline{\nabla u_2}\,\di x &= \redel \int_\Omega (A_1 - \I \AI_1)\nabla u_1\cdot\overline{\nabla u_2}\,\di x \notag \\
			&= \redel \int_\Omega A_2\nabla u_2\cdot\overline{\nabla u_2}\,\di x - \redel \I\int_\Omega \AI_1\nabla u_1\cdot\overline{\nabla u_2}\,\di x \notag \\
			&= \int_\Omega \AR_2 \nabla u_2\cdot\overline{\nabla u_2}\,\di x + \imdel \int_\Omega \AI_1 \nabla u_1\cdot\overline{\nabla u_2}\,\di x.
		\end{align}
		%
		By swapping the indices in \eqref{eq:mixed1}, and using that $\AR_2$ and $\AI_2$ are self-adjoint, we also have
		%
		\begin{equation} \label{eq:mixed2}
			\redel \int_\Omega \AR_2\nabla u_1\cdot\overline{\nabla u_2}\,\di x = \int_\Omega \AR_1 \nabla u_1\cdot\overline{\nabla u_1}\,\di x - \imdel \int_\Omega \AI_2 \nabla u_1\cdot\overline{\nabla u_2}\,\di x.
		\end{equation}
		
		Using \eqref{eq:mixed1} and that $\AR_1$ is positive definite, we have
		%
		\begin{align*}
			\int_\Omega (\AR_2 - \AR_1)\nabla u_2\cdot\overline{\nabla u_2}\,\di x + 2\imdel\int_\Omega \AI_1\nabla u_1\cdot\overline{\nabla u_2}\,\di x \hspace{-7cm}& \\
			&= \int_\Omega (\AR_2 - \AR_1)\nabla u_2\cdot\overline{\nabla u_2}\,\di x + 2\redel \int_\Omega \AR_1\nabla u_1\cdot\overline{\nabla u_2}\,\di x - 2\int_\Omega \AR_2 \nabla u_2\cdot\overline{\nabla u_2}\,\di x \\
			&\leq \int_\Omega (\AR_2 - \AR_1)\nabla u_2\cdot\overline{\nabla u_2}\,\di x + 2\redel \int_\Omega \AR_1\nabla u_1\cdot\overline{\nabla u_2}\,\di x - 2\int_\Omega \AR_2 \nabla u_2\cdot\overline{\nabla u_2}\,\di x\\
			&\phantom{={}} + \int_\Omega \AR_1\nabla(u_1-u_2)\cdot\overline{\nabla(u_1-u_2)}\,\di x\\
			&= \int_\Omega \AR_1\nabla u_1\cdot\overline{\nabla u_1}\,\di x - \int_\Omega \AR_2 \nabla u_2\cdot\overline{\nabla u_2}\,\di x.
		\end{align*}
		%
		That is, 
		%
		\begin{equation*}
			\inner{f,(\LambdaR_1-\LambdaR_2)f} \geq \int_\Omega (\AR_2 - \AR_1)\nabla u_2\cdot\overline{\nabla u_2}\,\di x + 2\imdel\int_\Omega \AI_1\nabla u_1\cdot\overline{\nabla u_2}\,\di x.
		\end{equation*}
		
		Using \eqref{eq:mixed2}, and that $\AR_2(\AR_1)^{-1}(\AR_2 - \AR_1)$ is self-adjoint, we obtain
		%
		\begin{align*}
			\int_\Omega \AR_2(\AR_1)^{-1}(\AR_2 - \AR_1) \nabla u_2\cdot\overline{\nabla u_2}\,\di x + 2\imdel \int_\Omega \AI_2\nabla u_1\cdot\overline{\nabla u_2}\,\di x \hspace{-9.5cm}& \\
			&= \int_\Omega \AR_2(\AR_1)^{-1}(\AR_2 - \AR_1) \nabla u_2\cdot\overline{\nabla u_2}\,\di x + 2\int_\Omega \AR_1 \nabla u_1\cdot\overline{\nabla u_1}\,\di x - 2\redel \int_\Omega \AR_2\nabla u_1\cdot\overline{\nabla u_2}\,\di x \\
			&\geq \int_\Omega \AR_2(\AR_1)^{-1}(\AR_2 - \AR_1) \nabla u_2\cdot\overline{\nabla u_2}\,\di x + 2\int_\Omega \AR_1 \nabla u_1\cdot\overline{\nabla u_1}\,\di x - 2\redel \int_\Omega \AR_2\nabla u_1\cdot\overline{\nabla u_2}\,\di x \\
			&\phantom{={}} - \int_\Omega \bigl\lvert (\AR_1)^{1/2}\nabla u_1 - (\AR_1)^{-1/2}\AR_2\nabla u_2 \bigr\rvert^2\,\di x \\
			&= \int_\Omega \AR_1\nabla u_1\cdot\overline{\nabla u_1}\,\di x - \int_\Omega \AR_2 \nabla u_2\cdot\overline{\nabla u_2}\,\di x.
		\end{align*}
		%
		That is, 
		%
		\begin{equation*}
			\inner{f,(\LambdaR_1-\LambdaR_2)f} \leq \int_\Omega \AR_2(\AR_1)^{-1}(\AR_2 - \AR_1) \nabla u_2\cdot\overline{\nabla u_2}\,\di x + 2\imdel \int_\Omega \AI_2\nabla u_1\cdot\overline{\nabla u_2}\,\di x. \qedhere
		\end{equation*}
	\end{proof}

\end{document}
